% !TEX program = pdflatex
%
% Signing the Instrument, Not the Transaction
% Preprint. Comments welcome.
%
% Build:  latexmk -pdf main.tex
%
\documentclass[11pt,a4paper]{article}

\usepackage[T1]{fontenc}
\usepackage[utf8]{inputenc}
\usepackage[margin=1in]{geometry}
\usepackage[protrusion=true,expansion=false]{microtype}
\usepackage{booktabs}
\usepackage{xcolor}
\usepackage[numbers,sort&compress]{natbib}
\usepackage{amsmath}
\usepackage{enumitem}
\usepackage[hidelinks,breaklinks]{hyperref}

\hypersetup{
  pdftitle={Signing the Instrument, Not the Transaction},
  pdfauthor={Chay Sengtha},
  pdfsubject={Layered trust failure in national QR payment systems},
}

\newcommand{\code}[1]{\texttt{\small #1}}

\title{\textbf{Signing the Instrument, Not the Transaction}\\[0.4em]
  \large Layered trust failure in national QR payment systems,\\
  with a reference implementation for Cambodia}

\author{Chay Sengtha\thanks{National University of Management, Phnom Penh.
Correspondence: \texttt{sengtha@gmail.com}.}}

\date{August 2026\\[0.5em]\normalsize\emph{Preprint. Comments welcome.}}

\begin{document}
\maketitle

\begin{abstract}
\noindent
Cambodia's KHQR, the national EMVCo-based payment QR operated on the National
Bank of Cambodia's Bakong system, has been adopted widely enough that the public
now treats \emph{any} QR code as a payment instrument and extends to it the
trust it holds in the central bank. Attackers exploit this in two ways that are
routinely conflated: forgery, where a code is not what it claims to be, and
deception, where a genuine code, correctly signed and payable to a correctly
registered account, is presented under a false pretext.

We specify and implement KH-SQR, a signature scheme for both payment codes and
document credentials, and we use it as evidence for a claim about where its own
assurance stops. Cryptography closes forgery completely and does not touch
deception at all, because a control enforced at one layer does not govern
meaning created at the layer above it. Signing the code secures the instrument;
it cannot secure the transaction, because the deception occurs one layer up, in
the person. This is the same shape as the central finding of our earlier work on
sovereign hybrid blockchains, where base-layer identity secured custody of an
asset but not exposure to it, because the derivative claim was minted above the
enforcement point.

The contribution is fourfold: a four-layer decomposition of QR payment trust
that locates precisely where cryptographic assurance ends; a working,
MIT-licensed reference implementation with a language-neutral conformance suite
of forty cases, thirty-one of them negative; an explicit account of what the
design does not achieve and why each residual risk is unreachable from the code
layer; and a constructive corollary --- that where a control cannot be enforced
below the harm, three moves remain --- raise the control above it, make the
consequence reversible, or reallocate the loss. We give the first concretely as
a categorical national prohibition on URL-bearing payment and official QR codes,
paired with a public rule that requires no judgement to apply, and we use Japan,
where the same QR code sits on a reversible stored-value substrate, to identify
finality rather than cryptography as the variable that governs exposure. We report two respects in which our own deployment and container fail to
meet the specification's stated requirements, and we treat the absence of
Cambodia-specific incident data as a finding rather than a gap to be filled with
vendor figures.
\end{abstract}

\medskip
\noindent\textbf{Keywords:} QR payments, authorised push payment fraud, digital
signatures, trust boundaries, EMVCo, Cambodia, layered security.

\section{Introduction}

A QR code is a picture of some bytes. Nothing about the picture indicates who
produced it, and nothing about the act of scanning it indicates why. Both facts
are obvious when stated and neither is available to a person holding a phone in
front of a printed square in a market.

Cambodia is an unusually clear case of what follows. KHQR, the national
merchant-presented payment standard operated on the National Bank of Cambodia's
Bakong system, has succeeded on a scale that is easy to understate: the NBC
reports 608.32 million Bakong transactions in 2024, and roughly 4.5 million
accounts registered for KHQR payments by the end of that year, a 36.4 per cent
increase over 2023 \citep{nbc2024annual}. Payment by QR is not a niche
behaviour in Cambodia; for a large part of the population it is simply what
paying is.

That success has a second-order effect which is the subject of this paper. When
a single visual form becomes synonymous with a single action, the form stops
carrying information. ``QR'' and ``payment'' have collapsed into one another in
public understanding, and with them ``QR'' and ``sanctioned by the central
bank''. The result is a divergence between two boundaries that a designer would
like to coincide:

\begin{itemize}[leftmargin=1.4em]
  \item the \textbf{trust boundary}: the set of artefacts whose provenance the
    NBC and its participants actually control and can vouch for;
  \item the \textbf{credulity boundary}: the set of artefacts to which the
    public extends the trust it holds in the NBC.
\end{itemize}

\noindent The second is much larger than the first, and it is the difference
between them that attackers operate in. A sticker over a merchant's code is
inside the gap. So is a printed code on a leaflet claiming to verify a document.
So, crucially, is a completely genuine KHQR presented with a false story
attached.

Signing the codes closes part of this gap and, we will argue, exactly the part
that is easiest to close and least often exploited. The rest is not a
cryptographic problem, and the central claim of this paper is that it cannot be
made into one:

\begin{quote}
\emph{A control enforced at one layer does not govern meaning created at the
layer above it.}
\end{quote}

\noindent Instantiated here: a signature proves who issued a code and cannot
address why the payer scanned it. We regard this as a general result about
layered systems rather than a local observation about Cambodia, and
Section~\ref{sec:general} argues the generalisation, including its relationship
to the custody/exposure finding in our earlier work on regulated DeFi
\citep{sengtha2026csb} and its applicability to other real-time push-payment
markets.

We are deliberate about what this paper is not. It is not a claim to have solved
QR phishing. A reviewer who finishes it believing that has been given the wrong
paper, and we have written Section~\ref{sec:notachieved} --- the longest
substantive section --- to make that reading difficult.

\paragraph{Contributions.}
\begin{enumerate}[leftmargin=1.6em]
  \item A four-layer decomposition of QR payment trust (Section~\ref{sec:layers})
    in which each layer's residual risk motivates the next, and in which the
    boundary between layers two and three is identified as the point where
    cryptography stops.
  \item KH-SQR: a specification and a working, MIT-licensed reference
    implementation covering a payment profile bound to EMVCo payloads and a
    standalone credential profile, with a language-neutral conformance suite of
    forty cases, thirty-one negative \citep{khsqr2026repo}.
  \item An explicit account of the residual risk (Section~\ref{sec:notachieved}),
    explaining for each item \emph{why} it is unreachable from the code layer
    rather than merely noting that it is, and an institutional-layer design
    (Section~\ref{sec:institutional}) that addresses part of it by
    non-cryptographic means.
  \item A policy instrument (Section~\ref{sec:categorical}) that reaches an
    attack the software cannot: a categorical prohibition on URL-bearing QR
    codes across the regulated perimeter, paired with an unconditional public
    rule. We show that every software alternative fails at the same point ---
    it requires a judgement before or during the scan --- and that the rule must
    not be scoped, because scoping it would require the user to classify a code
    before it is readable.
  \item Two honest negative results about our own work: the reference deployment
    does not satisfy the mirror-independence requirement the specification
    itself states, and the signature container is not transparent to a strict
    EMVCo parser as the design's premise assumed.
\end{enumerate}

\paragraph{On evidence.}
There is no reliable Cambodia-specific data on QR-code fraud. We searched for
it and did not find it, and we regard that absence as a finding worth stating
(Section~\ref{sec:data}) rather than a gap to be papered over with the vendor
figures currently in circulation. Where we make a technical claim, it is
supported either by a citation to a primary specification or by an artefact in
the repository that a reader can execute.

\section{Background}
\label{sec:background}

\subsection{KHQR and Bakong}

KHQR follows the EMVCo merchant-presented QR specification
\citep{emvco2017mpm}, which defines a payload as a flat sequence of
ID/Length/Value data objects. Both the ID and the Length are two decimal
digits; a Length ranges from ``01'' to ``99'', so no single data object may
carry more than 99 characters of value. The payload ends with a CRC-16 in tag
\code{63}. IDs \code{80} to \code{99} are ``unreserved templates'', available
for payment systems and value-added service providers to place their own data,
each required to carry a Globally Unique Identifier at sub-tag \code{00} that
establishes the template's namespace. These details matter to
Section~\ref{sec:design} and, uncomfortably for us, to
Section~\ref{sec:container}.

Bakong is the settlement layer beneath it. Its scale is documented in the NBC's
annual reporting \citep{nbc2024annual}; the property that matters for this
paper is not scale but \emph{finality}. A KHQR payment is a push payment: the
payer initiates and authorises a transfer to a payee, and once made it is
final. There is no chargeback rail comparable to card networks'.

\subsection{Irrevocability and the shape of the fraud economy}

Irrevocability is a design virtue with consequences on both sides of the ledger.
For merchants it removes chargeback risk and settlement delay, which is a
substantial part of why acceptance costs are low enough for a market stall. For
fraud it removes the recovery mechanism that card networks use to make deception
partially reversible.

The consequence is structural: in a push-payment system, an attacker who can
persuade a victim to authorise a payment has already won, and no property of the
payment instrument can undo it. Effort therefore migrates away from compromising
the instrument and towards compromising the persuasion. This is not a
prediction; it is the observed trajectory in every real-time push-payment market
with published statistics, and it is why the United Kingdom eventually addressed
its equivalent problem through liability allocation rather than instrument
security \citep{psr2024app}.

\subsection{The QR trust model, such as it is}

The trust model a user actually operates with, when scanning a code, is roughly:
\emph{this square was placed here by someone entitled to place it, and doing
what it says is what I intended to do}. Neither clause is verifiable by
inspection. A QR code is not human-readable; there is no equivalent of noticing
that a URL is misspelled, because there is no visible URL.

Existing mitigations largely attempt to move the decision into a channel the
user can inspect --- displaying the destination domain before opening it, showing
the payee name before confirming. These help, and Section~\ref{sec:notachieved}
gives the empirical reason to expect them to help less than designers hope.

\section{Threat decomposition}
\label{sec:layers}

We decompose the trust question into four layers. The decomposition is not
merely descriptive: each layer's residual risk is what motivates the next, and
the boundary between the second and third is the subject of the paper.

\begin{table}[t]
\centering
\small
\begin{tabular}{@{}p{0.13\linewidth}p{0.24\linewidth}p{0.26\linewidth}p{0.28\linewidth}@{}}
\toprule
\textbf{Layer} & \textbf{Question} & \textbf{Mechanism} & \textbf{Residual risk} \\
\midrule
1. Code authenticity &
  Was this code produced by a registered issuer, unaltered? &
  Digital signature over the payload; trust list; revocation &
  The verifier itself may be counterfeit \\[0.4em]
2. Verifier authenticity &
  Is the application making the claim genuine? &
  Application trust list; platform distribution; attestation &
  A genuine verifier faithfully reports a genuine code paid to a criminal \\[0.4em]
3. Payee legitimacy &
  \emph{Should} this money go to this account? &
  Registration controls; account risk lists; behavioural analytics &
  Legitimacy is a fact about intent and history, not about bytes \\[0.4em]
4. Recovery &
  What happens after a loss? &
  Liability allocation; reimbursement; freeze and recall procedures &
  Recovery is partial, slow, and after the fact \\
\bottomrule
\end{tabular}
\caption{Four layers of QR payment trust. Cryptography is decisive at layers~1
and~2 and silent at layers~3 and~4.}
\label{tab:layers}
\end{table}

\paragraph{Layer 1 --- code authenticity.}
Solvable, and solved by the design in Section~\ref{sec:design}. A signature over
the payload, verified against a trust list rooted in an offline key, establishes
that a registered issuer produced exactly these bytes. Within the standard
cryptographic threat model this is complete: an attacker cannot forge a
signature, and any modification invalidates one.

\paragraph{Layer 2 --- verifier authenticity.}
Partially solvable. A signature is only as useful as the software reporting it,
and a counterfeit wallet that draws a green tick for every code defeats layer~1
entirely without breaking any cryptography. Mitigations exist --- an application
trust list, platform distribution channels, device attestation --- and all of them
are weaker than layer~1's, because they depend on the user having obtained the
verifier through a trustworthy path.

\paragraph{The boundary.}
Between layers~2 and~3 the question changes kind. Layers~1 and~2 ask whether an
artefact is what it claims to be: a question about bytes, with a decidable
answer, computable offline from the artefact and a public key. Layer~3 asks
whether a payment \emph{ought} to be made: a question about intent, history and
context, whose answer is not a function of the artefact at all.

No amount of additional signing crosses this boundary, because there is nothing
on the far side to sign. A genuine merchant's code, correctly signed, remains
correctly signed when a criminal photographs it and pastes it onto a fake
utility-bill notice. Every byte is authentic. The fraud is entirely in the
sentence printed above the square.

\paragraph{Layer 3 --- payee legitimacy.}
Not solvable at the code layer, addressable in part at the institutional layer
by mechanisms that are emphatically not cryptography: registration controls,
account risk lists, behavioural detection. Section~\ref{sec:institutional}
specifies our contribution here.

\paragraph{Layer 4 --- recovery.}
Not a security control at all but an allocation of loss. Its design determines
who bears the cost of the residue the first three layers leave, and therefore
who has the incentive to reduce it.

\section{Design}
\label{sec:design}

KH-SQR defines two profiles. The full normative text is in the repository's
\code{SPEC.md} \citep{khsqr2026repo}; we describe here the decisions that carry
argumentative weight.

\subsection{Profile A: payment}

Profile A adds a signature to an existing EMVCo merchant-presented payload using
template \code{85} from the unreserved range, carrying a format version, a key
identifier, an algorithm identifier, an issuance time, an optional expiry, a
payee class, and a 128-character signature at sub-tag \code{99}.

\paragraph{The signing input is a prefix.}
The signed region is defined as the ASCII octets of the payload from position
zero up to \emph{and including} the five characters \code{99128} --- sub-tag
\code{99}'s tag and its length declaration --- excluding the signature itself and
the trailing CRC. Two ordering rules make this well defined: template \code{85}
must be the last data object before the CRC, and sub-tag \code{99} must be last
within it.

The consequence is that a verifier recovers the signed region with a substring
operation and never by re-serialising parsed fields. There is no canonical form
for two implementations to disagree about, and therefore no canonicalisation
step in which to have a bug. This is the single most important structural
decision in the design, and it is a deliberate rejection of the pattern --- parse,
normalise, re-serialise, verify --- that has produced a long history of signature
bypasses in XML, JSON and JWT processing.

\paragraph{Raw signatures, not DER.}
Signatures are the P1363 concatenation $r \Vert s$, 64 bytes, rendered as 128
uppercase hexadecimal characters. DER is prohibited. DER's length is variable,
which would defeat the fixed-offset property above, and DER length parsing is
itself a recurring source of defects. The choice also happens to eliminate a
conversion step entirely on the Web Crypto platform, whose ECDSA operations
consume and produce exactly this representation \citep{webcrypto,fips1865}.

\paragraph{Uppercase hexadecimal, not base64.}
QR alphanumeric mode encodes two characters in 11 bits and admits only digits,
uppercase letters, space and eight further punctuation marks
\citep{iso18004}. A single lowercase character forces a byte-mode segment at 8
bits per character. Base64's mixed case would therefore enlarge every signed
symbol.

\paragraph{Static and dynamic codes.}
A dynamic code carries an expiry no more than 300 seconds after issuance. A
static code carries no amount, so that a printed sticker cannot be given a
figure the merchant did not choose.

\subsection{Profile B: credential}

Profile B signs a document attestation as a standalone payload:
\[
\text{claims (CBOR)} \rightarrow \text{COSE\_Sign1 (ES256)} \rightarrow
\text{deflate (zlib)} \rightarrow \text{base45} \rightarrow \text{prefix
\code{KH1:}}
\]

The pipeline is the EU Digital COVID Certificate's, with a different prefix and
claim set \citep{eudcc2021,rfc9052,rfc9053,rfc8949,rfc1950,rfc1951,rfc9285}. We
say this plainly because the originality of this work is not in the primitives
or the pipeline; it is in the composition and in the analysis of what the
composition cannot do. Base45 rather than base64 for the same alphanumeric-mode
reason as above --- the choice the DCC made, for the same reason.

\paragraph{Excluding URL carriers.}
A Profile B payload must not be an \code{http} or \code{https} URL, and a
verifier must reject one. This is a design constraint, not an omission. A
URL-bearing QR relocates the trust decision into a browser and asks the user to
judge a domain name, which is the judgement that Section~\ref{sec:notachieved}
gives empirical reason to believe people do not reliably make. Our
implementation exposes this check as a function intended to be applied to
\emph{every} scanned code, not only to KH-SQR ones.

\paragraph{The transplant attack, enforced in the API shape.}
A valid signature proves a credential was issued. It does not prove the
credential belongs to the paper it is printed on. A genuine QR code photographed
from a real degree certificate and printed onto a forged one verifies perfectly,
because nothing about the paper is signed.

The only defence is comparison of the signed fields against the visible
document, and a defence that depends on a caller remembering to perform it will
not be performed. The specification therefore constrains the API rather than
merely advising: verification must return a structure carrying the subject name,
document identifier, issuing organisation and issue date in a form the caller
must handle; it must not return a boolean; and it must not expose an
\code{isValid} accessor by which a caller could reach a verdict without them. In
the reference implementation the field is named \code{mustMatchPrintedDocument},
and a test asserts that no boolean member exists on the result type --- so that a
later refactor adding a convenience accessor fails the build rather than
silently removing the only defence.

\subsection{Trust hierarchy and the fail-closed verifier}

Keys are identified by the first eight bytes of SHA-256 over the uncompressed
public point. The identifier is a lookup hint and never an authenticator: a
verifier encountering two trust-list entries with the same identifier must try
both and accept only on a signature that verifies.

The trust list is signed by an offline Root and carries a monotonic version and
an expiry. A verifier must reject a list numbered below the one it holds
(rollback), must stop verifying once its cached copy exceeds thirty days
(staleness), and must reject when the freshest \emph{timestamp statement} has
expired.

The timestamp statement is a separate signed object with a seven-day validity,
naming the current trust list's version and digest. It exists because expiry
alone does not defend against an attacker who can withhold updates: a hostile
network can pin a verifier to an old but still-valid list, so that a key revoked
yesterday keeps verifying. This is TUF's timestamp role, adopted for the same
reason TUF has it and cited as such \citep{tufspec}. The verifier fails closed:
absent a fresh timestamp statement it stops verifying rather than falling back
on what it holds.

Verification performs no network access. Everything it needs is supplied by the
caller. A verifier that fetched during verification could be stalled or steered
by whoever controls the network at the moment of payment, who at a market stall
is not a trustworthy party. The repository asserts this with a test that
poisons every network global and verifies both reference payloads successfully.

\subsection{No signing key at the edge}

The serving infrastructure holds no private key of any kind. The Root signs
offline in a ceremony; issuer keys live in each institution's hardware security
module; the timestamp statement is produced by a signer outside the serving
infrastructure and uploaded. The registry service accepts certificate signing
requests and queues them; it cannot issue.

This is specification rather than caution: certificate issuance must be
impossible through compromise of the online portal. Because the property is
invisible in the type system and easy to lose in a refactor, it is enforced by a
continuous-integration check that fails if a signing key appears in any service
configuration or source, and we verified that the check fails on a deliberately
introduced violation rather than merely passing on the current tree.

\section{Implementation and measurement}
\label{sec:implementation}

The reference implementation is TypeScript in strict mode, using the Web
Cryptography API exclusively \citep{webcrypto}, with the services deployed as
Cloudflare Workers. It is MIT licensed \citep{khsqr2026repo}.

\subsection{Dependency isolation}

Profile A's transitive import graph reaches six modules and depends on nothing
beyond Web Crypto: no CBOR, no stream APIs, no packages. The intent is
embeddability --- a mobile wallet should be able to take payment verification and
nothing else, with no polyfill. Profile B additionally uses CBOR and
\code{DecompressionStream}, which is acceptable because Profile B verifiers are
institutional (desktop and web) rather than consumer handsets. A build check
enforces the split and, like the key check, was verified to fail on an
introduced violation.

The CBOR codec is hand-written and covers only the fixed claim shapes in use. A
general CBOR library is a large audit surface on a security-critical path for no
benefit at this payload complexity. The decoder is deliberately strict ---
definite lengths only, minimal integer encodings only, no floating point, no
duplicate map keys, no trailing bytes --- because every construct accepted is a
construct an attacker may use to make two parsers disagree about the same bytes.
It is differentially fuzzed against an independent implementation in continuous
integration.

\subsection{Conformance suite}

The conformance suite is the artefact we expect to outlast the TypeScript. It is
a single JSON file of forty cases, thirty-one of them negative, each carrying an
identifier, a profile, a case type, an input, the trust state and time at which
it is evaluated, the expected outcome and, for a rejection, the expected
machine-readable reason. A Kotlin or Swift port can demonstrate conformance
without reading any of our code.

The negative cases are the substance. An implementation that accepts a
well-formed payload has demonstrated very little; one that rejects a
DER-encoded signature, a data object appended after the signature template, a
sub-tag appended after the signature, a byte mutated inside the signed region, a
byte mutated \emph{outside} it, an expired dynamic code, a static code carrying
an amount, an unknown key identifier, a revoked key, a rolled-back trust list, a
trust list past its cache age, an expired timestamp statement, an
\code{https} payload and a \code{deflate-raw} compression --- each for the correct
stated reason --- has demonstrated most of the specification.

Two case types are distinguished. A \emph{verify} case asserts that a recorded
payload produces a stated outcome. A \emph{roundtrip} case asserts that the
implementation can verify its own output. The distinction is necessary because
neither ECDSA nor DEFLATE is canonical: signing the same input twice yields
different signatures, and two conforming compressors produce different bytes. We
therefore require that the published Profile B payload \emph{decodes and
verifies}, not that an encoder reproduces it.

\subsection{Symbol size measurement}
\label{sec:measurement}

Table~\ref{tab:qr} reports measured QR symbol requirements at error-correction
level M. The measurements are generated by a script in the repository rather
than transcribed.

\begin{table}[t]
\centering
\small
\begin{tabular}{@{}lrrlp{0.3\linewidth}@{}}
\toprule
\textbf{Payload} & \textbf{Chars} & \textbf{Version} & \textbf{Modules} & \textbf{Encoding mode} \\
\midrule
Unsigned KHQR baseline        & 111 &  5 & $37\times37$ & numeric + byte + alphanumeric \\
Profile A signed              & 317 & 10 & $57\times57$ & numeric + byte + alphanumeric \\
Profile B credential          & 381 & 12 & $65\times65$ & alphanumeric \\
\addlinespace
Unsigned, uppercase acquirer id & 111 &  5 & $37\times37$ & numeric + alphanumeric \\
Profile A, uppercase acquirer id & 317 & 10 & $57\times57$ & numeric + alphanumeric \\
\bottomrule
\end{tabular}
\caption{Measured QR symbol requirements, error-correction level M
\citep{iso18004}. Generated by \code{tools/measure-qr.ts}.}
\label{tab:qr}
\end{table}

Signing takes a merchant payment code from version~5 to version~10: a 54 per
cent increase in linear dimension and 2.4 times the module count. The cost is
not abstract. At a fixed module size the printed sticker becomes larger; at a
fixed sticker size the modules become smaller and scan less reliably on an
inexpensive handset in poor light. It falls on merchants, who must reprint, and
a scheme that requires every merchant in the country to reprint has an adoption
problem that no amount of cryptographic elegance addresses.

The measurement also records the encoding mode, which surfaces something the
version alone conceals. The reference payload's acquirer identifier is
lowercase, which forces a byte-mode segment; making it uppercase keeps the
symbol alphanumeric throughout. At these lengths that does not change the
version, so the finding is a caution against attributing the size cost to the
wrong cause: the signature is what moves the symbol from version~5 to
version~10, not the case of one field.

Note that the EU DCC recommends error-correction level Q rather than M
\citep{eudcc2021}; our measurements at level M are therefore not directly
comparable to DCC symbol sizes, and a deployment choosing Q would see larger
symbols than Table~\ref{tab:qr} reports.

\subsection{The container is not legacy-transparent}
\label{sec:container}

The design's premise for using the unreserved template range was that existing
applications would ignore the added template while conforming ones required it.
\textbf{We report that this premise does not hold for a template of this size,
and that our container deviates from EMVCo in two respects.}

First, EMVCo length fields are exactly two decimal digits with a maximum value
of 99 \citep{emvco2017mpm}. Template \code{85} carries 201 characters and
sub-tag \code{99} carries 128. KH-SQR therefore declares both lengths in three
digits. A strict legacy EMVCo parser reading tag \code{85} would interpret the
first two of those digits as the length, consume 20 characters, misparse
everything following, fail to find the CRC tag, and reject the payload
outright --- rather than ignoring the template as intended.

Second, EMVCo requires that an unreserved template carry a Globally Unique
Identifier at sub-tag \code{00}, establishing the template's namespace
\citep{emvco2017mpm}. KH-SQR uses sub-tag \code{00} for a format version, so
template \code{85} does not identify itself in the way the specification
prescribes.

Neither deviation is a security defect: the length characters lie inside the
signed prefix, so tampering with them invalidates the signature, and our
verifier locates the signature by position rather than by any length field. But
the deployment consequence is real and was not in the original design's
reasoning. A national rollout must upgrade parsers, not merely signers, and the
migration story is correspondingly harder. A future revision should place a
genuine unique identifier at sub-tag \code{00} and move the format version
elsewhere; that is a wire-format change we have not made, because it would
invalidate the published reference vectors this work is cited against.

We also record an erratum in the published vectors: the reference payload
declares template \code{85}'s length as 200 where its content is 201 characters.
The payload is retained unchanged as a verification vector, since its signature
is valid and a conforming verifier accepts it; conforming signers emit the
correct value.

\subsection{Mirror independence is not satisfied}
\label{sec:mirrors}

The specification requires the trust list to be published at three mirrors under
\emph{distinct operational control}. Our reference deployment places all three
behind one provider.

Anycast delivers availability, not operational independence. Three URLs on one
provider share one account, one contractual relationship, one legal
jurisdiction, and one governance failure. If that account is suspended,
compromised, or subjected to a legal order, all three mirrors fail together, and
the verifier population enters its fail-closed staleness path at once.

\textbf{We therefore do not claim conformance to this clause}, and the service's
own health endpoint reports that it does not. Meeting it requires at least two
mirrors elsewhere --- the NBC's own infrastructure and a second commercial
provider are the obvious candidates. We state this because a specification that
requires operational independence and an implementation that quietly does not
provide it is precisely the kind of gap that a deployment inherits without
noticing.

\section{What the design does not achieve}
\label{sec:notachieved}

This is the most important section of the paper. For each item we give the
reason it is unreachable from the code layer, not merely the observation that it
is.

\subsection{Authorised push payment fraud}

A criminal shows a victim a genuine KHQR --- their own, correctly registered,
correctly signed --- with a false story: a refund to be released, a parcel to be
cleared, a fee to be paid. The victim scans, reads the payee name, and
authorises.

\textbf{Why it is unreachable.} Every byte of the payload is authentic, and a
correct implementation of this specification verifies it. There is no forgery to
detect because nothing was forged. The falsehood exists in the sentence spoken
or printed above the code, which is not part of the payload, is not available to
the verifier, and could not be signed by anyone even in principle --- there is no
party in a position to attest to the truth of an arbitrary claim about why a
payment is owed.

This is the layer boundary of Section~\ref{sec:layers} in its sharpest form. The
control is enforced on the instrument; the deception is created in the
interaction. One cannot govern the other because they are not the same kind of
object.

The design's only honest response is interface discipline: surfacing the payee
to the payer, and refusing to let a verification result be reduced to a tick.
Our specification makes both normative. We do not claim they are sufficient, and
Section~\ref{sec:indicators} explains why we expect them to be less effective
than they appear.

\subsection{Registration abuse}

A fraudster registers a merchant account using genuine documents --- their own, or
a coerced or purchased identity --- and receives a genuine key. The codes they
produce are correctly signed by a registered issuer, because that is exactly
what they are.

\textbf{Why it is unreachable.} The signature attests to an identity that is
real. Its truthfulness is not in question; its \emph{legitimacy} is, and
legitimacy is a fact about intent and history rather than about bytes. No
property computable from the payload distinguishes a genuine merchant from a
genuine criminal who completed the same registration process. The defence is
know-your-customer quality and post-hoc detection, both of which live at layer~3
and neither of which is cryptography.

\subsection{Counterfeit verifiers}
\label{sec:indicators}

An application that shows a green tick for every code defeats the entire scheme
without breaking any cryptography. An application trust list helps a diligent
user check what they installed and does nothing for a user who installed the
application from a link in a message.

\textbf{Why it is unreachable, and why security indicators disappoint.} The
relevant empirical result is Schechter et al.'s evaluation of website
authentication indicators \citep{schechter2007emperor}. In a study of 67 bank
customers, among the 25 participants who used their own accounts, 23 --- 92 per
cent --- entered their passwords after the site-authentication image, which they
had previously selected and were expected to verify, was removed and replaced
with an upgrade message.

We cite this precisely because it is the closest well-controlled evidence to the
question at hand, and we note its limits: it studied website images rather than
QR verification, in 2007, in a different population. It does not transfer
automatically to Cambodian merchants in 2026. What it establishes is narrower
and still decisive for design purposes: the \emph{absence} of a positive
security indicator is not reliably noticed by users who have been trained to
expect it. A scheme whose safety property is ``the user will notice when the
tick is missing'' rests on an assumption that has been tested and did not hold.

This is why the design excludes URL carriers rather than displaying domains
prominently, and why the credential API refuses to expose a boolean. Both are
attempts to avoid depending on a user noticing an absence.

\subsection{Foreign and platform codes}

Codes issued by systems outside the national trust hierarchy --- foreign wallets,
international platforms, ticketing and loyalty systems --- cannot be verified.

\textbf{Why it is unreachable, and why it is worse than it looks.} This is not
merely a coverage gap; it actively damages layers~1 and~2. A verifier that
displays ``unverified'' for the majority of codes a person encounters trains
that person to proceed past the warning. The indicator's value decays in
proportion to how often it fires benignly, which is the same mechanism that made
browser certificate warnings ineffective. A national scheme that verifies only
domestic codes may therefore leave users \emph{less} attentive than one that
makes no claim at all --- a possibility we flag rather than resolve, since we have
no data on it.

\subsection{The scan path: verification the user never invokes}
\label{sec:scanpath}

Every control in Section~\ref{sec:design} runs inside a conforming wallet. None
of it runs when the code is scanned with the handset's native camera
application, which is how a great many codes are in fact scanned --- a camera is
already open, it recognises the square, and it offers to act on it.

\textbf{Why it is unreachable.} The verifier cannot intervene in a scan it is
never asked to perform. The exclusion of URL carriers
(Section~\ref{sec:design}) is the clearest case: the rule is sound and the
implementation enforces it, but a URL-bearing code scanned by the camera
application opens in a browser without ever reaching any code we wrote. Closing
this in software would require the operating-system vendors to ship a
Cambodia-specific rule, which is not within national control and would not reach
the installed base of inexpensive handsets in any case.

This is the residual risk that motivates
Section~\ref{sec:categorical}, and it is why that section's intervention is not
software.

\subsection{Summary}

Layers~1 and~2 are the parts of the problem that cryptography addresses, and
this design addresses them. Layers~3 and~4 are the parts that matter most in
current Cambodian fraud, and it does not address them at all. Any presentation
of this work that inverts that emphasis misrepresents it.

There is, however, a constructive corollary that the layer argument itself
supplies, and it is the subject of Section~\ref{sec:categorical}. If the harm is
created above the enforcement point, the response is to raise the enforcement
point rather than to strengthen the control beneath it. The layer above the code
is not empty: it contains the rules governing what legitimate issuers are
permitted to emit. A control placed there does not authenticate the illegitimate
set; it empties the legitimate one, until anything outside becomes recognisable
without authentication.

\section{The institutional layer}
\label{sec:institutional}

Layers~3 and~4 are not reachable from the code, but they are not unreachable
altogether. They are reachable from the institutions, by two mechanisms that are
not cryptography: constraining what legitimate issuers are permitted to emit,
and sharing information about accounts. We take them in that order, because the
first is the cheaper and the more decisive.

\subsection{The categorical rule: a payment QR never opens a website}
\label{sec:categorical}

The URL-bearing QR code --- ``scan to claim your refund'', ``scan to verify your
document'' --- is the one attack in this decomposition that a national authority
can eliminate outright. It cannot be eliminated in software, for the reason
given in Section~\ref{sec:scanpath}: the code is frequently scanned by an
application that will never consult a trust list. It can be eliminated by rule.

\paragraph{The rule has two asymmetric halves.}
The asymmetry is the mechanism, and it is deliberate: the population that bears
the rule and the population that bears the obligation are different, and each is
given the form of instruction it can actually follow.

\begin{quote}
\textbf{Public rule.} A QR code never opens a website. If a website opens, it is
a scam: do not pay, and do not enter your PIN, password, one-time code or
personal details.
\end{quote}

\begin{quote}
\textbf{Institutional prohibition.} No licensed bank, payment institution,
government body or telecommunications operator shall issue, publish or display a
QR code whose payload is an \code{http} or \code{https} URL.
\end{quote}

The public half is categorical. It admits no exception, requires no comparison,
and is evaluated on an observation --- a browser window opened --- that anyone can
make regardless of literacy. The institutional half is a conditional compliance
obligation, which institutions can follow because they have compliance
functions and the public does not.

\paragraph{Why the rule must not be scoped.}
The obvious refinement --- restrict the prohibition to payments and official
communications, leaving restaurant menus and network logins alone --- fails, and
its failure is instructive. To apply a scoped rule, a person must classify the
code before acting on it; but a QR code is not human-readable, so its class is
knowable only after scanning, which is after the browser has opened. A scoped
rule is therefore unusable at the moment the decision is made. The categorical
rule is evaluated after the scan, on what is now visible, and requires no
classification at all. In a population with low general technology literacy this
is not a matter of preference: a rule with an exception is a rule that cannot be
applied.

Stating the rule in terms of the \emph{action} rather than the code --- do not
pay, do not authenticate --- keeps it categorical while leaving it uncontradicted
by the legitimate URL codes that exist outside the regulated perimeter. A
restaurant menu never asks for a one-time code, so encountering one does not
teach the public that the rule is false. This matters, because a public rule
falsified by daily experience decays exactly as the browser certificate warning
did.

\paragraph{Why there is no alternative that reaches the same attack.}
Every alternative we are aware of fails for one of two reasons: it requires a
judgement to be made \emph{before or during} the scan, which is precisely the
judgement the population cannot make; or it requires control over software the
state does not have.

\begin{table}[htbp]
\centering
\small
\begin{tabular}{@{}p{0.22\linewidth}p{0.26\linewidth}p{0.10\linewidth}p{0.32\linewidth}@{}}
\toprule
\textbf{Alternative} & \textbf{What it asks of the user} & \textbf{When} & \textbf{Why it fails} \\
\midrule
Wallet-only scanning &
  Choose the correct application &
  Before &
  The code's class is unknowable until it has been scanned \\[0.3em]
Display the destination domain &
  Read and compare a string &
  During &
  The judgement the population cannot reliably make \citep{schechter2007emperor} \\[0.3em]
In-wallet domain allowlist &
  Nothing &
  During &
  Bypassed entirely by the native camera application (Section~\ref{sec:scanpath}) \\[0.3em]
Operating-system camera integration &
  Nothing &
  During &
  Not within national control; unavailable on the installed handset base \\[0.3em]
Education without a categorical rule &
  Apply conditions &
  Before &
  Conditions are the thing that fails \\[0.3em]
Confirmation of payee &
  Check a name &
  After &
  Addresses the false pretext, not the URL channel; a different layer \\
\bottomrule
\end{tabular}
\caption{Alternatives to the categorical rule, and the decision point at which
each one fails.}
\label{tab:alternatives}
\end{table}

The categorical rule is the only intervention in
Table~\ref{tab:alternatives} evaluated \emph{after} the scan, on an observation
requiring no comparison and no prior classification. That is not an assertion
that nothing better exists; it is a property of where each control sits relative
to the decision point, which is the same argument the rest of this paper runs on.

\paragraph{Who is obliged, and why.}
The public trust that makes this attack profitable is not a natural fact. The
identification of ``QR'' with ``payment'', and of ``payment'' with
``sanctioned by the central bank'', is the output of a successful national
campaign --- it is what adoption at the scale reported in Section~\ref{sec:background}
means. A trust expectation manufactured by policy can only be bounded by policy,
and by the same authority that created it. This is a statement about standing,
not about blame: no other party can make the rule true, because no other party
can bind the institutions whose conduct would otherwise falsify it.

Accordingly the requirement is stated flatly rather than conditionally. Whether
it is adopted is a political question and not a property of the design.

\begin{itemize}[leftmargin=1.4em]
  \item The prohibition MUST bind the full regulated perimeter --- payment
    institutions, government bodies and telecommunications operators alike.
    A single ministry running a URL-QR campaign falsifies the public rule for
    everyone, and the public cannot be asked to hold an exception.
  \item Legacy printed material MUST be withdrawn within a stated sunset period,
    with an audit. Branch posters, ATM decals and existing marketing are the
    counter-examples most likely to be encountered.
  \item A named regulator MUST hold the obligation, with a public channel for
    reporting violations by institutions.
  \item The public wording MUST be identical across institutions. Variation is
    how a categorical rule erodes into a conditional one.
\end{itemize}

Section~3.2 of the specification is the enforcement half of the same rule inside
the software: a Profile~B payload must not be an \code{http} or \code{https}
URL, and the reference implementation exposes that check as a function intended
to be applied to every scanned code rather than only to KH-SQR ones
\citep{khsqr2026repo}. Neither half is sufficient alone. The software half is
bypassed by the scan path; the policy half has no technical enforcement against
an attacker, only against the legitimate issuers whose compliance makes it true.

\paragraph{What the rule does not do.}
It closes a channel; it does not reduce the adversary's budget. Effort displaces
to what remains permitted --- most obviously to the genuine-code-with-false-pretext
attack of Section~\ref{sec:notachieved}, which is already the dominant vector and
which this rule leaves entirely intact --- and to channels with no QR code at all.
We claim the elimination of one attack, not a reduction in fraud, and the
distinction is the same one this paper insists on everywhere else.

\subsection{Time-to-list as the operative metric}

A mule account's useful life is short: it receives, it forwards, it is
abandoned. The quantity that determines how much passes through it is the
interval between the first institution recognising it and every other
institution being able to act --- \emph{time-to-list}. Detection accuracy matters
less than dissemination latency, because an accurate signal delivered after the
funds have moved has no effect on the loss.

This has an immediate architectural consequence, and it is the reason our risk
list is not built on an eventually consistent store. An eventually consistent
read leaves a window in which a just-listed account still reads clear, and that
window is exactly when the account is being drained. The implementation
therefore routes reads and writes for each account through a serialising object,
with a durable database as the authority.

\subsection{Screening: the checkpoint the list exists for}
\label{sec:screening}

A register that nobody consults before releasing money is a record of the fraud,
not a control on it. The operative step is therefore a check made by the payer's
institution, at the moment of payment, against the payee account named in the
code.

This is also the one point at which layer~1 feeds layer~3. A verified Profile~A
attestation yields the payee account identifier from a payload that has been
proved unaltered, so the account screened is the account that will be paid. We
should be precise about how much the signature contributes here: the account
identifier would be present in an unsigned payload too, and the risk list would
still work on it. What the signature adds is that the identifier cannot have
been substituted in transit, and that the code's issuer is attributable
afterwards. The control that actually bites on this fraud is the list, not the
signature.

The reference implementation exposes screening as a single call returning a
decision rather than a status, because a status invites each institution to
invent its own mapping, and a national scheme in which one bank holds where
another releases is not a national scheme. The mapping is normative:

\begin{center}
\small
\begin{tabular}{@{}lll@{}}
\toprule
\textbf{Status} & \textbf{Action} & \textbf{Character} \\
\midrule
\code{clear}      & Execute & --- \\
\code{restricted} & Hold: delay, or route to manual review & Prudential, reversible, expires in 72 hours \\
\code{blocked}    & Refuse & Standing assertion, two officers \\
\bottomrule
\end{tabular}
\end{center}

A low-value carve-out is available --- below a threshold, a payment to a
restricted account may proceed with a prominent warning rather than being held
--- because holding a one-dollar transfer costs a real person more inconvenience
than it denies an attacker. It is disabled by default in every currency. The
safe value is the default, and a regulator who wants less friction must choose
the threshold and own the choice.

\paragraph{A hold is borrowed reversibility.}
The deeper justification for holding rather than merely warning is not caution.
A system whose payments are final has, by construction, no way to undo a
transfer once the payer has been deceived --- and that property, not the quality
of anyone's cryptography, is the variable Section~\ref{sec:japan} identifies as
separating Cambodia's exposure from Japan's. A hold reintroduces a small,
targeted quantity of the missing reversibility: for the narrow set of payments
that trip a signal, and for a bounded period, the transfer is not yet final. It
does not restore general reversibility, which Cambodia gave up deliberately and
for good reasons. It buys back just enough of it, in just the cases where the
loss would otherwise be irreversible.

\paragraph{What is recorded, and what deliberately is not.}
Only decisions with a consequence --- warn, hold, block --- are written down. A
cleared payment leaves no trace, and the payer is never identified to the
service at all. Recording every screening would build a national record of who
paid whom, which is not this system's business and is not needed for its
purpose: whether an account was listed at a given moment is already
reconstructable from the append-only change feed. We note this because a
transaction-screening service is exactly the kind of infrastructure that
acquires surveillance functions by accident, one reasonable-sounding field at a
time.

\paragraph{What screening does not do.}
It stops payments to accounts that are \emph{already listed}. The victims of a
mule account during the interval before anyone lists it receive nothing, which
is why time-to-list and not detection accuracy is the metric. It sees the first
hop only, and nothing of cash-out or onward layering. And it sits in the payment
path, so if it is not fast, institutions will find ways around it.

\subsection{A hold is not a sanction}

The design distinguishes two statuses, and the distinction is a matter of
principle rather than granularity:

\begin{itemize}[leftmargin=1.4em]
  \item \textbf{Restricted} --- a prudential hold. One institution's operational
    judgement about its own exposure, taking effect immediately, expiring by
    default after seventy-two hours unless renewed.
  \item \textbf{Blocked} --- a standing assertion about an account, requiring
    approval by a second officer before it takes effect.
\end{itemize}

Conflating them would let an operational reflex acquire the force of a penalty.
A hold that never expires is a sanction imposed without process; the expiry is
what keeps it a hold. Removal likewise requires two officers, in both
directions, so that a single compromised or coerced credential can neither
freeze money nor release it.

\subsection{Expiry is evaluated when the status is read}

Statuses carry a stored deadline compared against the clock when someone asks,
and there is no sweep job. This is a deliberate choice about failure direction.
A cron that fails to run would silently extend a restriction on a real person's
account, and neither they nor the institution would observe anything wrong; the
error is invisible and prejudicial. A deadline that has simply passed cannot
fail in that direction.

\subsection{The right to contest a listing}
\label{sec:appeal}

A national register that can freeze a real person's money needs a way for that
person to be wrong about. Expiry, two-officer approval and an immutable audit
trail all reduce the chance of error. None of them lets the affected person do
anything about one.

The design has three constraints, and the third is the substantive one.

First, the customer cannot query this service and cannot appeal to it directly.
An open lookup would tell a mule operator whether their account had been
detected yet, which would convert a fraud control into reconnaissance. The
account-holding institution raises the contest on its customer's behalf.

Second, raising a contest does not change the status. An institution must not be
able to clear a suspicion --- its own or another's --- by asserting that it is
disputed.

Third, \textbf{an unanswered contest lapses the listing.} Raising it starts a
clock against the institution that made the listing: 24 hours for a restriction,
72 hours for a block, in both cases shorter than the listing's own lifetime so
that appealing accelerates something real. If that institution does not answer
before the deadline, the listing ceases to be in force --- evaluated at read time,
like the expiry, with no sweep job. Silence favours the account holder, because
the account holder is the party who cannot act. An institution that wishes to
keep an account listed must be willing to say so, on the record, with a named
officer against it.

\paragraph{An asymmetry we had to correct.}
Building this exposed a defect in the design as originally specified. A
restriction takes one officer to impose but two to remove, which makes an error
more expensive to correct than to make --- the wrong way round for a control that
acts on people who have done nothing. The rule is now narrower: a withdrawal by
the institution that made the listing, in answer to a contest, takes one officer,
because it retracts that institution's own assertion and creates no new one. The
two-officer requirement stays where it belongs, on a discretionary removal by an
institution that did not make the listing. Upholding likewise takes one officer,
so that answering a contest is never more onerous than ignoring it.

\paragraph{An unresolved tension.}
The listing institution's identity is disclosed to the account-holding
institution, not to the customer. This preserves the ability to contest while
respecting the tipping-off constraints that anti-money-laundering regimes
impose, but it means the affected person learns that they are listed and by
which category of reason, without learning who listed them or on what evidence.
We do not think this is satisfactory, and we do not have a better answer that
does not compromise investigations. We record it as an open problem rather than
resolving it by assertion.

\subsection{Attribution, rate limiting and the audit record}

Every write records the institution \emph{and} the individual officer, via
per-officer mutual TLS credentials. An audit entry naming only an institution is
not an audit entry: accountability requires that someone specific be nameable
afterwards. The log is append-only, enforced by database triggers rather than by
convention, and hash-chained so it can be exported as a file whose interior
integrity is checkable without trusting the database it came from, and published
to a transparency log later without redesign. Corrections are new rows.

Writes are rate limited per institution and per officer, and sustained volume
above a threshold is refused \emph{and} recorded as an incident rather than
merely throttled. An institution listing thousands of accounts in an hour is
compromised or misconfigured; neither condition is a reason to act on its
assertions.

\subsection{Liability allocation}
\label{sec:liability}

The fourth layer is not a security control but an allocation of loss, and its
design determines who has the incentive to reduce the residue.

The United Kingdom's approach is the clearest available precedent. Since 7
October 2024, in-scope authorised push payment scam victims must be reimbursed
up to a maximum of GBP 85,000, with the cost shared equally between the sending
and the receiving payment service provider \citep{psr2024app}. The equal split is
the substantive design choice: it gives the receiving institution --- the one that
opened the mule account and is best placed to detect it --- a direct financial
interest in the quality of its own onboarding, which a sender-pays rule would
not.

We note the precedent rather than recommending its transplantation. Cambodia's
market structure, dollarisation, and the balance between banks and payment
institutions all differ, and a reimbursement cap calibrated to UK incomes has no
meaning elsewhere. What transfers is the principle that liability should rest
where the cheapest available control is, and that a scheme which leaves the
entire loss with the payer has by construction allocated it to the party with
the least ability to prevent it.

We have not implemented anything at this layer. It is policy, and it is named
here because a threat decomposition that stopped at layer~3 would misrepresent
where the residual risk finally settles.

\section{Generalisation}
\label{sec:general}

\subsection{The layer result}

Stated generally: \emph{a control enforced at layer $n$ cannot govern semantics
created at layer $n+1$.} A signature binds bytes to an issuer. If the harm is
produced by the meaning a person attaches to those bytes, and that meaning is
created above the signed layer, then no strengthening of the signature reduces
the harm. Effort spent there is not merely insufficient; it is misdirected, and
it carries the additional cost of producing an assurance that invites the
inference it does not support.

The diagnostic question for any proposed control is therefore not ``is this
mechanism sound?'' but ``at which layer is the harmful meaning created, and is
that at or below the layer where this mechanism is enforced?''

\subsection{Three degrees of freedom}
\label{sec:degrees}

If a control cannot reach the layer at which the harm is created, three moves
remain, and we can identify only three. Each acts at a different point in the
causal chain: before the meaning is created, on the consequence of acting on it,
or on who bears that consequence.

\begin{table}[htbp]
\centering
\small
\begin{tabular}{@{}p{0.19\linewidth}p{0.36\linewidth}p{0.37\linewidth}@{}}
\toprule
\textbf{Move} & \textbf{Mechanism} & \textbf{Instance in this paper} \\
\midrule
Raise the control above the harm &
  Empty the set of legitimate artefacts, so anything outside it is recognisable
  without authentication &
  The categorical prohibition on URL-bearing codes
  (Section~\ref{sec:categorical}) \\[0.4em]
Make the consequence reversible &
  Interpose a ledger somebody controls, so a transfer can be undone after the
  deception has already succeeded &
  Japan's stored-value substrate (Section~\ref{sec:japan}); and, narrowly, the
  prudential hold (Section~\ref{sec:screening}) \\[0.4em]
Reallocate the loss &
  Move the cost to the party holding the cheapest available control &
  Liability allocation (Section~\ref{sec:liability}) \\
\bottomrule
\end{tabular}
\caption{What remains available when a control cannot reach the layer where the
harm is created.}
\label{tab:degrees}
\end{table}

The second move is the one worth dwelling on, because it is the least obvious
and because it does not contradict the layer result --- it sidesteps it.
Reversibility does not govern meaning at all. It governs \emph{consequence}. The
deception still succeeds, the victim is still deceived, and the harm is undone
from \emph{below} the layer at which the lie was told. That is possible
precisely because consequence, unlike meaning, is a property of the substrate,
and the substrate is something a designer chooses.

\subsection{Japan: the same code, a different substrate}
\label{sec:japan}

The QR code was invented in Japan, at Denso Wave in 1994, by a team led by
Masahiro Hara, to track automotive parts; its designers did not exercise their
patent rights, which is the main reason the format became universal
\citep{densowave}. Japan also has substantial code-payment volume
\citep{meti2026cashless}. It does not have Cambodia's QR fraud shape, and the
reason is not better cryptography.

Japanese code payments are predominantly operator-mediated stored value: the
scan spends a balance funded in advance, or authorises a linked card, and
settlement to the merchant occurs afterwards across the operator's own books.
Cambodia's arrangement is the opposite. The scan \emph{is} the settlement
instruction, and the transfer is interbank, immediate and final.

Four consequences follow from that single difference:

\begin{itemize}[leftmargin=1.4em]
  \item an operator is interposed, holding a ledger it controls and can adjust;
    in an interbank push there is no middle, and nobody positioned to reverse;
  \item the loss is bounded by the funded balance rather than by the account
    balance;
  \item recourse exists --- an operator dispute for prepaid value, and inherited
    chargeback rights where a card funds the wallet;
  \item the set of payees is curated by the operator, so paying an arbitrary
    stranger's account is not the default primitive.
\end{itemize}

\noindent Two qualifications, and both matter.

\textbf{Japan is not free of the deception this paper is about.} Impersonation
and investment fraud are serious there. The money travels by bank transfer, cash
handover or crypto instead. Stored value did not stop deception; it stopped
deception arriving \emph{by QR}. Attackers go where finality is, which is an
uncomfortable way to read Cambodia's exposure rather than a reassuring one.

\textbf{Finality was chosen deliberately, and it bought something.} Removing
chargeback risk is part of why acceptance costs are low enough for a market
stall (Section~\ref{sec:background}). Recommending Japan's architecture would
undo that and insert intermediary rent in its place. We do not recommend it. We
use it to identify the variable, and then --- in Section~\ref{sec:screening} --- to
buy a little of that variable back where it matters most.

We characterise the architecture at the level of its class rather than auditing
individual operators, and we report no Japanese incident figures: the sources we
consulted disagreed on both category and magnitude, and we did not verify them
against the National Police Agency's own releases.

\subsection{The parallel with custody and exposure}

Our earlier work on sovereign hybrid blockchains for regulated DeFi reached the
same shape by a different route \citep{sengtha2026csb}. Base-layer identity
controls secured \emph{custody} of an asset --- who may hold and move it --- and did
not secure \emph{exposure} to it, because the derivative claim conveying that
exposure was minted one layer above the enforcement point. The identity control
was correct and did what it said; the harm simply happened somewhere it did not
reach.

The correspondence is exact:

\begin{center}
\small
\begin{tabular}{@{}lll@{}}
\toprule
 & \textbf{Regulated DeFi} \citep{sengtha2026csb} & \textbf{QR payments (this work)} \\
\midrule
Control        & Base-layer identity   & Signature over the code \\
What it secures & Custody of the asset  & Authenticity of the instrument \\
What it misses  & Exposure to the asset & The transaction, and its pretext \\
Why             & Claim minted one layer above & Meaning created one layer above \\
\bottomrule
\end{tabular}
\end{center}

Two independent domains producing the same failure shape is weak evidence for
generality --- two cases is two cases --- but it is enough to motivate the
diagnostic question above as a design practice.

\subsection{Other push-payment markets}

Where Japan is a contrast, India's UPI and Brazil's PIX are instances. Both are
real-time push-payment systems in which payments are irrevocable and initiated
from a machine-readable artefact --- the same substrate as Bakong --- and both have
contended with the same migration of attacker effort from instrument to
persuasion. Contrast and instance do different work: the instances show the
pattern recurs, and the contrast identifies which property drives it. We do not claim
detailed knowledge of either market's incident data and make no quantitative
comparison. The structural point stands independently of the numbers: where
payments are final and initiation is delegated to an artefact the user cannot
read, securing the artefact is necessary and is not close to sufficient.

\section{Data availability and the absence of Cambodian incident statistics}
\label{sec:data}

We are unable to report the prevalence of QR-code fraud in Cambodia, and we
consider this worth stating explicitly.

There is no published national incident dataset. The figures currently in
circulation --- a widely syndicated claim of a 146 per cent quarter-on-quarter
increase in ``quishing'' in the first quarter of 2026, and similar --- originate in
vendor marketing material. They describe telemetry from the vendors' own
customer bases, with undisclosed denominators, no stated methodology, and no
Cambodian breakdown. They are not evidence for any claim in this paper and we do
not rely on them; we mention them only because a reader will encounter them and
should know why they are absent here.

The absence itself is a finding. A payment system at the scale the NBC reports
\citep{nbc2024annual} without public incident statistics cannot support
evidence-based prioritisation: no party can demonstrate that forgery is rarer
than deception, which is the very claim this paper argues on structural grounds
and would prefer to argue on data. Establishing mandatory incident reporting for
payment institutions, with published aggregates, would do more for the design of
countermeasures than any cryptographic improvement to the codes --- including
this one.

Everything empirical in this paper is therefore either a citation to a primary
specification or an artefact that a reader can execute
\citep{khsqr2026repo}. The symbol measurements in Table~\ref{tab:qr} are
produced by a script; the conformance claims are produced by a test suite; the
architectural properties are asserted by continuous-integration checks that we
verified fail on deliberately introduced violations.

\section{Limitations}

\begin{enumerate}[leftmargin=1.6em]
  \item \textbf{No deployment, no field evaluation.} Nothing here has been
    deployed with real merchants or real users. We report no usability data, no
    scan-reliability data at the larger symbol sizes, and no adoption data.
    Section~\ref{sec:measurement}'s claim that a larger symbol scans less well on
    an inexpensive handset is a reasonable inference, not a measurement.

  \item \textbf{Mirror independence is unmet.} Section~\ref{sec:mirrors}. The
    implementation does not satisfy a requirement its own specification states.

  \item \textbf{The container is not legacy-transparent.}
    Section~\ref{sec:container}. Two deviations from EMVCo, one of which
    invalidates the design's stated premise for the choice of template range.

  \item \textbf{The Japanese comparison is architectural, not empirical.}
    Section~\ref{sec:japan} characterises a class of payment architecture and
    reports no Japanese incident data. We did not audit individual operators,
    and we did not verify fraud figures against the National Police Agency's own
    releases. The contrast identifies a variable; it does not measure its effect.

  \item \textbf{The categorical rule is proposed, not adopted or evaluated.}
    Section~\ref{sec:categorical} argues from the structure of the decision
    point, not from evidence that the rule changes behaviour. We know of no
    trial of a categorical QR rule in any jurisdiction, and we have not run one.
    Its cost --- the legitimate URL-QR uses it forecloses inside the regulated
    perimeter --- is real and unquantified here.

  \item \textbf{Screening and appeals are implemented but unevaluated.} The
    endpoints, the status-to-action mapping and the lapse-on-silence rule are
    built and tested (Sections~\ref{sec:screening} and~\ref{sec:appeal}), but we
    have no evidence about false-positive rates, about how often institutions
    would answer a contest in time, or about the latency screening would add in
    a live payment path.

  \item \textbf{Not implemented.} The offline Root ceremony is specified but not
    performed. Liability allocation (layer~4) is discussed and not implemented,
    being policy rather than software. Device attestation for verifier
    authenticity is named as a mitigation and not built. Cross-border trust-list
    federation is out of scope entirely.

  \item \textbf{Truncated key identifiers.} Eight bytes is a size choice, not a
    collision-resistance claim. The specification requires verifiers to treat
    the identifier as a hint and to try every candidate, which contains the
    consequence, but an implementer who ignores that rule has a real bug.

  \item \textbf{One implementation.} The conformance suite exists to enable
    independent ports and has not yet been run against one. Until it has, the
    suite tests our reading of our own specification.

  \item \textbf{No formal verification, no external audit.} The parsers are
    hand-written and differentially fuzzed, which is materially weaker than
    either.

  \item \textbf{The SiteKey evidence is analogical.}
    Section~\ref{sec:indicators}. A 2007 study of website images in a
    non-Cambodian population supports a narrow claim about unnoticed absent
    indicators; it does not establish how Cambodian users respond to QR
    verification interfaces, and we know of no study that does.
\end{enumerate}

\section{Conclusion}

We set out to build a signature scheme for Cambodia's national QR payments and
to say precisely what it achieves. It achieves the elimination of forgery: a
code that verifies was produced by a registered issuer and has not been altered,
and the reference implementation demonstrates this against a published
conformance suite whose negative cases outnumber its positive ones.

It achieves nothing whatever against the fraud that currently dominates. A
genuine code, correctly signed and payable to a correctly registered account,
presented with a false story, verifies --- and must. The deception is created one
layer above the enforcement point, and no control at the lower layer reaches it.
That is not a deficiency of this design but a property of layered systems, and
the same property we found in a different domain when base-layer identity
secured custody of an asset and left exposure to it untouched.

The practical implication is uncomfortable for anyone hoping cryptography will
settle the matter. The next marginal improvement in Cambodian QR fraud outcomes
will not come from a better signature. It will come from a categorical rule that
empties the legitimate set of URL-bearing payment codes, which costs nothing to
enforce against the small regulated population that would otherwise falsify it;
from shortening the interval between an institution recognising a mule account
and every other institution being able to act on it; from allocating liability
so that the party best placed to prevent a loss bears part of it; and, before
any of them, from publishing incident data so that the prioritisation can be
argued from evidence rather than, as here, from structure.

Each of those is an intervention at a layer the code cannot reach, which is the
point. The response to a control that stops too low is not a stronger control at
the same height.

We offer the artefact because the artefact is useful and the forgery layer
genuinely needs closing. We offer the analysis because the artefact will be
misread as doing more than it does, and the most valuable thing a specification
can contain is an accurate statement of its own limits.

\section*{Declarations}

\paragraph{Funding.} This work received no external funding.

\paragraph{Competing interests.} The author declares no competing interests.

\paragraph{Data and code availability.} The reference implementation,
conformance vectors, measurement scripts and specification are available at
\url{https://github.com/sengtha/qrseal} under the MIT licence, cited at tag
\texttt{v1.0.0} \citep{khsqr2026repo}. The published test key material is
deliberately public and protects nothing. All tables and quantitative claims in
this paper are reproducible from that repository.

\paragraph{Disclaimer.} This is a preprint and has not been peer reviewed. It
describes a research design and a reference implementation. It is not an
endorsement by, nor produced in cooperation with, the National Bank of Cambodia
or any operator of the Bakong system, and nothing in it should be read as a
statement of their plans or positions. The implementation has not been
independently audited and is not represented as fit for production deployment.
Nothing here constitutes legal or regulatory advice.

\paragraph{Generative AI disclosure.} Generative AI tools were used in the
preparation of this work: for drafting and editing prose, for implementing and
testing the reference code, and for locating candidate sources. All technical
claims were verified by the author against primary specifications or by
executing the accompanying artefacts; the test vectors were independently
re-derived and confirmed; and the design decisions, the layer analysis and the
argument are the author's. The author accepts full responsibility for the
content, including any errors.

\bibliographystyle{plainnat}
\bibliography{references}

\end{document}
